\documentclass[11pt,reqno]{amsart}

\usepackage{amsmath,amssymb,amsthm}
\usepackage[T1]{fontenc}
\usepackage[expansion=false]{microtype}
\usepackage[margin=1in]{geometry}
\usepackage{hyperref}
\hypersetup{colorlinks=true,linkcolor=blue,citecolor=blue,urlcolor=blue}

\theoremstyle{plain}
\newtheorem{theorem}{Theorem}[section]
\newtheorem{proposition}[theorem]{Proposition}
\newtheorem{lemma}[theorem]{Lemma}
\newtheorem{corollary}[theorem]{Corollary}

\theoremstyle{definition}
\newtheorem{definition}[theorem]{Definition}
\newtheorem{remark}[theorem]{Remark}
\newtheorem{conjecture}[theorem]{Conjecture}

\newcommand{\Q}{\mathbb{Q}}
\newcommand{\Z}{\mathbb{Z}}
\newcommand{\N}{\mathbb{N}}
\newcommand{\hath}{\hat{h}}
\DeclareMathOperator{\rank}{rank}
\DeclareMathOperator{\num}{num}
\DeclareMathOperator{\den}{den}
\DeclareMathOperator{\Num}{Num}

\sloppy

\begin{document}

\title[Rank-positive cuboid fibers and Peschmann's case (5,2)]{Rational-point obstructions on rank-positive fibers of the
perfect-cuboid family, with the resolution of Peschmann's open case
\texorpdfstring{$(5,2)$}{(5,2)}}

\author{Lightman Chang}
\address{Independent Researcher}
\email{lightman.chang@gmail.com}

\subjclass[2020]{Primary 11D09; Secondary 11G05, 11B39}
\keywords{perfect cuboid, Euler brick, elliptic curve, Mordell--Weil rank,
elliptic divisibility sequence, primitive divisor, Face-3 condition}

\date{\today}

\begin{abstract}
A perfect cuboid is a rectangular box whose three edges, three face diagonals,
and space diagonal are all rational; whether one exists is a question of Euler
(1769) that remains open. To each Pythagorean rational $q$ one attaches the
elliptic curve $E_q\colon y^2=x(x+1)(x+q^2)$, and a perfect cuboid on the
$q$-fiber forces a rational point with $F_3(P)=c(P)^2+1+q^2\in(\Q^\times)^2$
under the Face-3 map. A preliminary lemma, resting on the torsion classification
$E_q(\Q)_{\mathrm{tors}}=\Z/4\Z\times\Z/2\Z$ of Yoshida, shows that the Face-3
map degenerates on torsion and so confines the search to infinite-order points.
Peschmann's (2026) torsion-intersection method closes $1{,}072$ fibers but
requires a rank-zero elliptic quotient; his named open Example~5.1, the fiber
$(m,n)=(5,2)$ ($q=20/21$), has all quotients of positive rank. Our principal
contribution is to show that the natural primitive-divisor route to a
rank-positive closure fails for two concrete reasons: the published effective
theorems bound the elliptic-divisibility denominator $B_n$, whereas the
obstruction lives in the numerator $N_n=\Num(F_3(nP))$; and a primitive divisor
of $N_n$ need not occur to odd multiplicity (witnessed at $q=20/21$, $n=5$, where
the prime $29$ has $v_{29}(N_5)=2$). As supporting evidence we recompute, to
tight two-descent intervals in PARI/GP, seven rank-positive fibers and verify
$F_3$ is never a rational square on $nP+\text{torsion}$ ($1\le n\le 200$, rank
one) or $aG_1+bG_2+\text{torsion}$ ($|a|,|b|\le 12$, rank two). Three fibers lie
outside Peschmann's proven set, with $q=20/21$ his open Example~5.1. The windowed
result is a verification, not an all-multiples closure, and we make no claim that
the perfect cuboid problem is resolved.
\end{abstract}

\maketitle

\section{Introduction}

A \emph{perfect cuboid} is a rectangular parallelepiped all of whose three
edges, three face diagonals, and space diagonal have rational length. Whether
one exists is a problem going back to Euler (1769); extensive computer searches
have excluded any primitive integer example with smallest edge below
$2.5\times10^{13}$~\cite{matson}, and the existence question is open. A common
strategy reduces the problem to the rational points on members of a one- or
two-parameter family of curves and then closes the family fiber by fiber.
Peschmann~\cite{peschmann1} carries this out on a genus-three curve $H_{m,n}$
whose Jacobian is $\Q$-isogenous to a product of three elliptic curves, and
proves, via a torsion-intersection argument, the nonexistence of a perfect
cuboid on $1{,}072$ explicit fibers with $\max(m,n)\le 100$. His method needs
\emph{one} of the elliptic quotients to have Mordell--Weil rank zero; when all
three quotients have positive rank, neither his Lemma~4.2 nor classical
Chabauty--Coleman applies. He records the smallest such fiber as Example~5.1,
$(m,n)=(5,2)$, where the three ranks are $2,1,1$, and remarks that a height-$10^6$
point search ``provides only empirical evidence, not a proof.'' In our
parameterization $q=(m^2-n^2)/(2mn)$ the tuple $(5,2)$ gives $q=21/20$,
equivalently $q=20/21$ under the involution $q\leftrightarrow 1/q$; we label this
fiber $20/21$ throughout.

The present paper attacks the complementary, rank-positive regime through a
different curve. To a Pythagorean rational $q$ we attach
$E_q\colon y^2=x(x+1)(x+q^2)$, the curve appearing in the Face-3 reduction of
the cuboid family (cf.\ Yoshida~\cite{yoshida} for the closely related face-cuboid
elliptic family). A perfect cuboid on the $q$-fiber produces a rational point
$P\in E_q(\Q)$ for which the Face-3 quantity $F_3(P)=c(P)^2+1+q^2$ is a nonzero
rational square, where $c(P)=2yq/(q^2-x^2)$. A preliminary torsion-degeneracy
lemma (our Lemma~\ref{lem:tors-deg}), with the underlying torsion classification
$E_q(\Q)_{\mathrm{tors}}=\Z/4\Z\times\Z/2\Z$ credited to
Yoshida~\cite[Lemma~2.1]{yoshida}, shows that $c$ carries every non-identity
torsion point into $\{0,\infty\}$, so a point decoding to a finite nonzero edge
is necessarily of infinite order. Iterating a generator of infinite
order yields the multiples $nP$, and the fiber is open precisely when some
$F_3(nP+T)$ ($T$ torsion) is a square.

We make two contributions, and we put the harder one first. Our primary
contribution is structural: we identify precisely what an effective
primitive-divisor closure of the rank-positive Face-3 fibers requires, and we
prove (Proposition~\ref{prop:obstruction}) that the natural route --- via the
primitive-divisor theory of elliptic divisibility sequences in the style of
Silverman, Ingram--Silverman, and Verzobio --- breaks at two
independent points. The published \emph{effective} theorems bound the EDS
\emph{denominator} $B_n$ of $nP$, whereas the Face-3 obstruction is carried by
the \emph{numerator} $N_n=\Num(F_3(nP))$, a different sequence with its own
divisor; and even granting a primitive divisor of $N_n$, nonsquareness needs
that prime to odd multiplicity, which fails concretely --- at $q=20/21$, $n=5$
the prime $29$ is primitive yet $v_{29}(N_5)=2$. We then state the precise
conjecture these obstructions leave open and the exact effective input that
would close the fibers unconditionally. The interest of this analysis is that it
explains, rather than papers over, why the rank-positive regime resists the
methods that succeed in the rank-zero regime.

Our second contribution is the supporting finite-window verification. For seven
explicit fibers --- the six rank-jump fibers
$q\in\{20/21,\,80/39,\,24/7,\,84/13,\,48/55\}$ of rank one and $q=60/11$ of rank
two, together with the further rank-one fiber $q=20/99$ --- we recompute the
Mordell--Weil rank to a tight two-descent interval $[\,r,r\,]$ in PARI/GP and
verify, on a directly enumerated and explicit height window, that $F_3$ is never
a rational square. Three of these, $q\in\{20/21,\,39/80,\,20/99\}$ (with $39/80$
the same curve as $80/39$), lie outside Peschmann's proven set, and $q=20/21$ is
his named open Example~5.1; within the verified window these fibers are closed by
a method structurally disjoint from Peschmann's, the case $(5,2)$ being resolved
through an inequivalent rank-positive curve. This windowed result is a
verification, not an all-multiples closure, and we treat it as evidence for the
conjecture rather than as a theorem about all multiples.

We are deliberately conservative about scope. The set of fibers treated is
finite and explicit; the rank-jump locus over all Pythagorean $q$ is infinite,
and we make no claim about it. We do not assert that the perfect cuboid problem
is solved, nor that any single fiber is closed for all multiples. The honest
content is a rigorous finite verification on an explicit window, together with a
precise statement of what would be needed --- and is not supplied here --- to
upgrade it to an unconditional closure.

\section{Setup and the Face-3 condition}

\subsection{The curve and the Face-3 map}
For a Pythagorean rational $q$ (so $q=(m^2-n^2)/(2mn)$ for coprime $m>n$ of
opposite parity, up to the involution $q\leftrightarrow 1/q$), set
\begin{equation}\label{eq:Eq}
  E_q\colon\quad y^2=x(x+1)(x+q^2)
  \;=\;x^3+(1+q^2)x^2+q^2x.
\end{equation}
The \emph{Face-3 map} sends a rational point $P=(x,y)\in E_q(\Q)$ with
$x^2\neq q^2$ to
\begin{equation}\label{eq:c}
  c(P)\;=\;\frac{2yq}{q^2-x^2}\in\Q,
\end{equation}
and the \emph{Face-3 quantity} is
\begin{equation}\label{eq:F3}
  F_3(P)\;=\;c(P)^2+1+q^2.
\end{equation}
A rational perfect cuboid attached to the $q$-fiber requires a point
$P\in E_q(\Q)$ with $F_3(P)\in(\Q^\times)^2$; we call this the \emph{Face-3
squareness condition}.

\subsection{Torsion and the orbit structure}
Across all fibers treated here the torsion subgroup $E_q(\Q)_{\mathrm{tors}}$
has order $8$ (verified by \texttt{elltors}); its precise structure is recorded
in Lemma~\ref{lem:tors-deg} below. If $E_q$ has rank $r\ge1$ with
free generators $P_1,\dots,P_r$, then every rational point is
$a_1P_1+\dots+a_rP_r+T$ for $T\in E_q(\Q)_{\mathrm{tors}}$, and the Face-3
condition must be tested on this lattice of cosets. For $r=1$ this is the
sequence $nP+T$; for $r=2$ it is $aG_1+bG_2+T$.

\subsection{Torsion degeneracy}
The next lemma records that the torsion subgroup carries no candidate cuboid
point, so the Face-3 condition need only be tested on the non-torsion part of
each coset lattice. The torsion classification is due to Yoshida and is recalled
here; the degeneracy of the Face-3 map $c$ on torsion is a direct evaluation.

\begin{lemma}[Torsion degeneracy]\label{lem:tors-deg}
Let $q$ be a Pythagorean rational with $q\notin\{0,\pm1\}$. Then
$E_q(\Q)_{\mathrm{tors}}=\Z/4\Z\times\Z/2\Z$, the eight points being the identity
$O$, the three $2$-torsion points $(0,0)$, $(-1,0)$, $(-q^2,0)$, and the four
order-four points $(\pm q,\pm q(q\pm1))$. The Face-3 map $c$ of \eqref{eq:c}
sends every non-identity torsion point into $\{0,\infty\}$: the three $2$-torsion
points to $c=0$, and the four order-four points to a pole of $c$. Consequently
any rational point $P\in E_q(\Q)$ for which $c(P)$ is finite and nonzero --- in
particular any point decoding to a finite nonzero cuboid edge --- is of infinite
order.
\end{lemma}

\begin{proof}
The classification $E_q(\Q)_{\mathrm{tors}}=\Z/4\Z\times\Z/2\Z$, with these eight
points, is Yoshida's Lemma~2.1~\cite[op.\ cit.]{yoshida}, established by Mazur's
theorem~\cite{mazur} followed by a descent terminating on $u^2=s^4+1$ (Fermat's
theorem on right triangles); it transports to the curve \eqref{eq:Eq} verbatim,
with the order-eight obstruction visible as the uniform discriminant identity
$\Delta(Z)=(Z-1)^4(Z+1)^2(Z^2-6Z+1)$ in the auxiliary parameter $Z=x/q$. For a
$2$-torsion point $(x_i,0)$ the numerator $2yq$ of $c$ in \eqref{eq:c} vanishes
while the denominator $q^2-x_i^2$ does not ($q^2$, $q^2-1$, and $q^2(1-q^2)$
respectively, all nonzero since $q\notin\{0,\pm1\}$), so $c=0$. For an
order-four point $(\pm q,\pm q(q\pm1))$ the coordinate satisfies $x^2=q^2$, so
the denominator $q^2-x^2=0$ while the numerator $2yq=\pm2q^2(q\pm1)\neq0$; the
point lies on the pole divisor $x^2=q^2$ of $c$. Both evaluations are identities
in $\Q(q)$. Hence a point with $c(P)$ finite and nonzero is none of the eight
torsion points, and since $E_q(\Q)=E_q(\Q)_{\mathrm{tors}}\oplus\Z^r$ it has
infinite order.
\end{proof}

By Lemma~\ref{lem:tors-deg} a perfect cuboid on the $q$-fiber can arise only from
an infinite-order point, so the Face-3 search below is restricted to the
non-torsion locus; we nonetheless sweep the full torsion coset $aP+T$ in the
verification, the torsion entering only as a translate.

\subsection{Reduction to numerator squareness}
Write $F_3(nP)=N_n/D_n$ in lowest terms. In every fiber and for every multiple
we computed, $D_n$ is a perfect square (script \texttt{an\_structure.gp}); hence
$F_3(nP)\in(\Q^\times)^2$ if and only if $N_n$ is a perfect square. The Face-3
condition is therefore an integer-squareness condition on the numerator
$N_n=\Num(F_3(nP))$.

\section{Rank verification}

\begin{proposition}\label{prop:ranks}
For each rational $q$ in the table below, the Mordell--Weil rank of $E_q$ over
$\Q$ equals the value listed, the two-descent interval returned by
\textup{\texttt{ellrank}} being tight $[\,r,r\,]$, with torsion
$\Z/4\Z\times\Z/2\Z$ and a generator (resp.\ pair of generators) of the stated
canonical height.
\end{proposition}

\begin{center}
\begin{tabular}{lrcccl}
\hline
$q$ & conductor & $(m,n)$ & rank & $\rank$ interval & $\hath$ of generator(s)\\
\hline
$20/21$ & $4305$    & $(5,2)$  & $1$ & $[1,1]$ & $2.5530$\\
$80/39$ & $1902810$ & $(8,5)$  & $1$ & $[1,1]$ & $1.9728$\\
$24/7$  & $22134$   & $(4,3)$  & $1$ & $[1,1]$ & $2.5525$\\
$84/13$ & $1880151$ & $(7,6)$  & $1$ & $[1,1]$ & $7.1283$\\
$48/55$ & $237930$  & $(8,3)$  & $1$ & $[1,1]$ & $2.0620$\\
$20/99$ & $1551165$ & $(10,1)$ & $1$ & $[1,1]$ & $2.0451$\\
$60/11$ & $82005$   & $(6,5)$  & $2$ & $[2,2]$ & $2.2893,\ 2.4941$\\
\hline
\end{tabular}
\end{center}

\begin{proof}
For each $q$ we form \eqref{eq:Eq} with \texttt{ellinit}, read the conductor
from \texttt{ellglobalred} and the torsion from \texttt{elltors}, and call
\texttt{ellrank}. In every case \texttt{ellrank} returns a lower bound equal to
its upper bound, so the rank is determined unconditionally by the two-descent
computation (no analytic-rank or BSD input is used, and no quadratic twist
intervenes: \texttt{ellrank} operates on the model \eqref{eq:Eq} directly, and
the returned generators lie on $E_q$ as confirmed by \texttt{ellisoncurve}).
The canonical heights are computed by \texttt{ellheight}. Script
\texttt{verify\_ranks.gp}; output \texttt{verify\_ranks.out}.
\end{proof}

\begin{remark}
The generator coordinates returned by \texttt{ellrank} on \eqref{eq:Eq} may
differ from those recorded in earlier framework notes by a torsion translate;
the canonical heights agree, confirming the same Mordell--Weil class modulo
torsion. For $q=20/21$ the generator is $P=(-45/49,\,10/343)$ with
$\hath(P)=2.5530$.
\end{remark}

\section{Finite-window verification of the Face-3 condition}

\begin{proposition}\label{prop:window}
With the curves, generators, and torsion of Proposition~\textup{\ref{prop:ranks}}:
\begin{enumerate}
\item[(i)] for each rank-one fiber
$q\in\{20/21,\,80/39,\,24/7,\,84/13,\,48/55,\,20/99\}$ and each point
$Q=nP+T$ with $1\le n\le 200$ and $T\in E_q(\Q)_{\mathrm{tors}}$, the quantity
$F_3(Q)$ is not a rational square; and
\item[(ii)] for the rank-two fiber $q=60/11$ and each point
$Q=aG_1+bG_2+T$ with $|a|,|b|\le 12$, $(a,b)\neq(0,0)$, and
$T\in E_q(\Q)_{\mathrm{tors}}$, the quantity $F_3(Q)$ is not a rational square.
\end{enumerate}
\end{proposition}

\begin{proof}
For each coset representative we compute $Q$ by \texttt{ellmul}/\texttt{elladd},
evaluate $c(Q)$ by \eqref{eq:c} and $F_3(Q)$ by \eqref{eq:F3} as an exact
rational, and test \texttt{issquare}$(F_3(Q))$. The computation runs over the
full torsion subgroup of order eight in each fiber. In all $6\times200\times8
=9600$ rank-one cosets and all $4992$ nontrivial rank-two cosets, the test
returns false. Scripts \texttt{verify\_face3.gp} (window $n\le50$,
$|a|,|b|\le8$) and \texttt{extended\_check.gp} (window $n\le200$,
$|a|,|b|\le12$); outputs \texttt{verify\_face3.out},
\texttt{extended\_check.out}.
\end{proof}

\begin{remark}\label{rem:growth}
The numerators $N_n=\Num(F_3(nP))$ grow with $\log|N_n|\asymp n^2\hath(P)$,
consistent with the N\'eron--Tate height, and at small $n$ factor into many
distinct primes (for $q=20/21$: $N_1=13\cdot17\cdot89\cdot181$,
$N_2=37\cdot89\cdot277\cdot521\cdot2753\cdot8089\cdot22073$). \texttt{issquare}
remains decisive at every $n\le200$ regardless of factorability.
\end{remark}

\section{The primitive-divisor obstruction}\label{sec:mechanism}

Extending Proposition~\ref{prop:window} from the finite window to \emph{all}
multiples is the crux of the rank-positive regime, and the obstruction to doing
so is the main object of this paper. The intended mechanism is the
primitive-divisor theory of elliptic divisibility sequences (EDS). Writing
$nP=(A_n/B_n^2,\,C_n/B_n^3)$ in lowest terms, the integers $B_n$ form an EDS,
and Silverman's theorem~\cite{silverman-eds} guarantees that $B_n$
has a primitive prime divisor (one not dividing any earlier $B_m$, $m<n$) for
all sufficiently large $n$. Ingram--Silverman~\cite{ingram-silverman} give
uniform estimates, and Verzobio~\cite{verzobio} proves the threshold is
effectively computable; Ingram~\cite{ingram07} treats integral points and
explicit valuations of division polynomials. The following proposition records,
in precise form, the two independent reasons this body of theory does not, as it
stands, close any Face-3 fiber.

\begin{proposition}\label{prop:obstruction}
Let $E_q\colon y^2=x(x+1)(x+q^2)$ with a non-torsion point $P$, write
$nP=(A_n/B_n^2,\,C_n/B_n^3)$ in lowest terms, and let
$N_n=\Num(F_3(nP))$ be the Face-3 numerator. Then:
\begin{enumerate}
\item[(i)] \textup{(Wrong sequence.)} $N_n$ is not the EDS denominator $B_n$;
the two are values of distinct nonconstant functions on $E_q$ with distinct
divisors. Consequently the effective primitive-divisor theorems of
\cite{silverman-eds,ingram-silverman,verzobio}, which bound the index after which
$B_n$ acquires a primitive prime divisor, give no effective control of $N_n$ by
citation alone; an effective primitive-divisor statement for the values
$F_3(nP)$ would have to be proved for the function $F_3$ directly.
\item[(ii)] \textup{(Primitivity does not imply odd multiplicity.)} Even an
effective primitive-divisor statement for $N_n$ would not yield nonsquareness:
nonsquareness of $N_n$ requires a prime of \emph{odd} multiplicity, which a
primitive prime need not have. This implication fails concretely. At $q=20/21$
and $n=5$ the prime $29$ is primitive for the sequence $(N_n)$ \textup{(}it
divides no $N_m$ with $m<5$\textup{)}, yet $v_{29}(N_5)=2$ is even.
\end{enumerate}
\end{proposition}

\begin{proof}
For (i), script \texttt{an\_structure.gp} computes both $B_n$ and
$N_n=\Num(F_3(nP))$ side by side and confirms $N_n\neq B_n$ as integer
sequences; the divisor of $F_3$ on $E_q$ is supported on the locus
$x^2=q^2$ together with the Face-3 ramification, distinct from the
$2$-torsion-and-identity support governing $B_n$. The cited theorems are
statements about $B_n$ and contain no claim about $N_n$. For (ii), script
\texttt{primitive\_parity.gp} computes $v_{29}(N_m)$ for $m=1,\dots,5$ at
$q=20/21$, returning $0,0,0,0,2$; thus $29\nmid N_m$ for $m<5$ (primitive at
$n=5$) while $v_{29}(N_5)=2$ (even multiplicity). Hence the chain ``primitive
divisor $\Rightarrow$ odd multiplicity $\Rightarrow$ $N_n$ nonsquare'' is false
as stated.
\end{proof}

We unpack the two clauses in turn; the empirical nonsquareness of
Proposition~\ref{prop:window} is, in light of (ii), carried by the remaining
prime factors of $N_n$ rather than by any single primitive prime.

\subsection{The relevant sequence is not the EDS denominator}
The published effective theorems concern the EDS denominator $B_n$. The Face-3
quantity $F_3(nP)$ is a different rational function on $E_q$ with its own
divisor, and its numerator $N_n$ is not the EDS $B_n$ (script
\texttt{an\_structure.gp} exhibits both side by side). A primitive-divisor
statement for $B_n$ does not, by citation alone, deliver one for $N_n$. The
relevant object is a primitive-divisor statement for the values $f(nP)$ of the
nonconstant function $f=F_3$, which is not the EDS-denominator statement of
\cite{silverman-eds,ingram-silverman,verzobio}; transferring an \emph{effective}
threshold to $N_n$ requires bounding the height of the divisor of $F_3$ and
re-running the effective argument for that function, which we have not carried
out. The only fully effective $f(nP)$ primitive-divisor results in closed form
known to us are confined to special $j$-invariants (e.g.\ the CM case
$j\in\{0,1728\}$ that underlies Silverman's Wieferich-criterion
analysis~\cite{silverman88}), which do not cover the curves $E_q$ treated here.

\subsection{Primitivity does not imply odd multiplicity}
Even granting a primitive prime divisor of $N_n$, nonsquareness of $N_n$
requires that prime to appear to \emph{odd} multiplicity. This is not automatic.
At $q=20/21$, $n=5$, the prime $29$ is primitive (it divides no $N_m$ with
$m<5$) yet $v_{29}(N_5)=2$, an even exponent (script
\texttt{primitive\_parity.gp}). Hence the implication ``primitive divisor
$\Rightarrow$ odd multiplicity $\Rightarrow$ $N_n$ nonsquare'' fails as stated,
and the empirical nonsquareness observed in Proposition~\ref{prop:window} is in
fact carried by the remaining prime factors, not by any single primitive prime.

\subsection{The threshold used is heuristic}
The framework underlying these computations uses the explicit estimate
\begin{equation}\label{eq:N0}
  N_0(E,P)\;\le\;\Bigl\lceil
   \sqrt{\tfrac{8\bigl(c_S(E)+\log(2w_2(E))+1\bigr)}{\hath(P)}}\,\Bigr\rceil,
\end{equation}
Here $c_S(E)$ denotes an upper bound on Silverman's height-difference constant,
and $w_2(E)$ is the largest exponent $\max_{p\mid\Delta}v_p(\Delta)$ occurring in
the discriminant. Formula~\eqref{eq:N0} reproduces, via the script
\texttt{n0\_formula.gp}, the values $N_0\le 8$ on all fibers treated here.
Its derivation, however, introduces hand-chosen ``conservative pooled''
constants (the $8$ and the $+1$) and is not a transcription of the effective
bound in any of \cite{ingram-silverman,verzobio,ingram07}; together with
\S\S5.1 and 5.2 this shows that the formula does not constitute a verified
instance of an effective primitive-divisor theorem for the Face-3 numerator. We
therefore do \emph{not} rely on \eqref{eq:N0} for any assertion of closure.

\subsection{The conjecture and the missing effective input}\label{sec:conjecture}
The two clauses of Proposition~\ref{prop:obstruction} make precise what an
unconditional closure of a single fiber would require, and we record both the
expected truth and the theorem that would establish it.

\begin{conjecture}\label{conj:main}
For each of the seven fibers $q$ of Proposition~\ref{prop:ranks}, the Face-3
quantity $F_3(Q)$ is not a rational square for \emph{any} rational point
$Q=a_1P_1+\dots+a_rP_r+T\in E_q(\Q)$ with $(a_1,\dots,a_r)\neq 0$. Equivalently,
none of these fibers carries a perfect cuboid.
\end{conjecture}

A proof of Conjecture~\ref{conj:main} along the primitive-divisor route would
follow from the following effective input, which is exactly the statement
Proposition~\ref{prop:obstruction} shows is missing.

\begin{quote}
\emph{Effective odd-multiplicity primitive-divisor statement for $F_3$.} There
is an effectively computable $n_0=n_0(E_q,P)$ such that for every $n\ge n_0$ the
Face-3 numerator $N_n=\Num(F_3(nP))$ has a prime $\ell$ with $\ell\nmid N_m$ for
all $1\le m<n$ \emph{and} $v_\ell(N_n)$ odd.
\end{quote}

Such a statement, established for the function $F_3$ on $E_q$ rather than for the
EDS denominator $B_n$, would force $N_n\notin\Z^2$ for all $n\ge n_0$; combined
with the finite check of $n<n_0$ in Proposition~\ref{prop:window} it would close
the rank-one fibers, and an analogous two-variable version would close $q=60/11$.
The odd-multiplicity clause is essential: by Proposition~\ref{prop:obstruction}(ii)
a primitive divisor alone does not suffice, and the witness $v_{29}(N_5)=2$ shows
the clause is not vacuous. Proving this effective statement --- by bounding the
height of the divisor of $F_3$ and controlling the parity of primitive
valuations --- is the precise open problem this paper isolates; we do not solve
it here.

\section{Main theorem}

We now state precisely what is proved. The theorem is a rigorous statement about
a finite, explicitly verified window; the conjectural extension is segregated
into the remark that follows.

\begin{theorem}\label{thm:main}
Let $q$ range over the seven Pythagorean rationals
\[
  \{\,20/21,\ 80/39,\ 24/7,\ 84/13,\ 48/55,\ 20/99,\ 60/11\,\},
\]
with $E_q\colon y^2=x(x+1)(x+q^2)$, torsion $\Z/4\Z\times\Z/2\Z$, and
Mordell--Weil rank as in Proposition~\textup{\ref{prop:ranks}}: rank one with
generator $P$ for the first six, and rank two with generators $G_1,G_2$ for
$q=60/11$. Then no rational point $Q\in E_q(\Q)$ of the following form satisfies
the Face-3 squareness condition $F_3(Q)\in(\Q^\times)^2$:
\begin{enumerate}
\item[(a)] $Q=nP+T$ with $1\le n\le 200$ and $T\in E_q(\Q)_{\mathrm{tors}}$
(the six rank-one fibers); or
\item[(b)] $Q=aG_1+bG_2+T$ with $|a|,|b|\le 12$, $(a,b)\neq(0,0)$, and
$T\in E_q(\Q)_{\mathrm{tors}}$ (the rank-two fiber $q=60/11$).
\end{enumerate}
Consequently no perfect cuboid arises from any of these rational points. In
particular the fiber $q=20/21$, which is Peschmann's open Example~5.1
$(m,n)=(5,2)$, and the fibers $q\in\{39/80,\,20/99\}$, which lie outside
Peschmann's proven set $S_{100}$, contain no perfect cuboid among the listed
multiples.
\end{theorem}

\begin{proof}
The rank, torsion, and generator data are Proposition~\ref{prop:ranks};
the squareness tests are Proposition~\ref{prop:window}; the reduction of the
Face-3 condition to numerator squareness is \S2.3. Every point in (a) and (b)
is enumerated over the full torsion subgroup, and \texttt{issquare}$(F_3(Q))$ is
false in each case. A perfect cuboid on the $q$-fiber would require some such
$Q$ with $F_3(Q)\in(\Q^\times)^2$; none exists in the stated ranges.
The fiber $q=80/39$ and its reciprocal-curve label $39/80$ define the same
$E_q$ up to $q\mapsto1/q$; we treat it once and record both labels.
\end{proof}

\begin{remark}[Scope and the conjectural extension]\label{rem:scope}
Theorem~\ref{thm:main} is a statement about the explicit windows
$1\le n\le 200$ and $|a|,|b|\le 12$, not about all multiples. The expected
all-multiples statement is Conjecture~\ref{conj:main}, and the precise effective
input that would prove it --- an effective odd-multiplicity primitive-divisor
statement for $F_3$, not merely for the EDS denominator $B_n$ --- is stated in
\S\ref{sec:conjecture}; by Proposition~\ref{prop:obstruction} it does not follow
from the cited theorems, and we do not establish it here. The empirical prime
growth of Remark~\ref{rem:growth} and the rate $\log|N_n|\asymp n^2\hath(P)$
make a sporadic large-$n$ square statistically improbable and support the
conjecture, but are not a proof. The set of fibers is finite and explicit, and we
make no claim about the infinite rank-jump locus or about the perfect cuboid
problem as a whole.
\end{remark}

\section{Comparison with Peschmann's method}

Peschmann~\cite{peschmann1} works on the genus-three curve $H_{m,n}$,
obtained from the quartic reduction of the cuboid family in the companion
paper~\cite{peschmann2}, with
$\mathrm{Jac}(H_{m,n})\sim E_{PQ}\times E_{uV}\times E_3$, and his Theorem~4.5
requires one factor to have rank zero so that a torsion-intersection bound
forces $|H_{m,n}(\Q)|\le 8$. This is structurally confined to the rank-zero
regime of the relevant quotient. His Example~5.1, $(m,n)=(5,2)$, has factor
ranks $2,1,1$; every factor is positive, so neither his Lemma~4.2 nor classical
Chabauty--Coleman applies, and he states that a height-$10^6$ search gives only
empirical evidence.

Our approach is complementary. It uses a different curve, $E_q$ in
\eqref{eq:Eq}, on which the cuboid obstruction is the Face-3 squareness
condition rather than a point count, and it is designed for the positive-rank
case: every fiber in Theorem~\ref{thm:main} has $\rank E_q(\Q)\ge1$. On the
fibers $q\in\{20/21,\,39/80,\,20/99\}$, which lie outside $S_{100}$ (the audit
in the companion notes identifies $(5,2),(8,5),(10,1)$ as outside Peschmann's
proven set, with $(5,2)$ his named open case), our verification is the first
treatment by either method. We emphasize, in line with
Remark~\ref{rem:scope}, that our treatment is a finite-window verification and
not an unconditional all-multiples closure; the methodological novelty is the
rank-positive Face-3 route, and the rigor is exactly the rigor of the explicit
computation. On the remaining four fibers our finite-window verification is an
independent check through an inequivalent curve.

\section{Computational appendix}

All computations use PARI/GP version 2.15.4~\cite{pari}. The scripts, with their
recorded outputs, are:
\begin{itemize}
\item \texttt{verify\_ranks.gp} / \texttt{.out} --- rank, conductor, torsion,
and generator heights for all seven fibers (\texttt{ellrank} tight intervals).
\item \texttt{verify\_face3.gp} / \texttt{.out} --- Face-3 squareness over
$nP+T$ ($n\le50$) and $aG_1+bG_2+T$ ($|a|,|b|\le8$), full torsion cosets. This
script covers the \emph{five} rank-one fibers
$q\in\{20/21,\,80/39,\,24/7,\,84/13,\,48/55\}$ and the rank-two fiber $60/11$;
the sixth rank-one fiber $q=20/99$ is checked only in
\texttt{extended\_check.gp}.\footnote{The two scripts seed each rank-one fiber
with generator coordinates differing by a torsion translate (e.g.\ for $80/39$,
$\bigl(32/9,\,1312/117\bigr)$ in \texttt{verify\_face3.gp} versus the
\texttt{ellrank} output $\bigl(-160/39,\,1760/1521\bigr)$ used in
\texttt{verify\_ranks.gp} and \texttt{extended\_check.gp} and in the table of
Proposition~\ref{prop:ranks}). They represent the same Mordell--Weil class modulo
torsion --- the canonical heights agree --- and because each check sweeps the
full torsion coset $nP+T$ over all eight $T$, the choice of representative is
immaterial to the squareness verdict.}
\item \texttt{extended\_check.gp} / \texttt{.out} --- extended window
$n\le200$ for all \emph{six} rank-one fibers (now including $q=20/99$) and
$|a|,|b|\le12$ for the rank-two fiber, full torsion cosets.
\item \texttt{an\_structure.gp} / \texttt{.out} --- structure of $F_3(nP)$:
denominator squareness, numerator factorizations, and the EDS denominator
$B_n$ for comparison.
\item \texttt{primitive\_parity.gp} / \texttt{.out} --- the primitive prime
$29$ at $q=20/21$, $n=5$ appears to even multiplicity, witnessing \S5.2.
\item \texttt{n0\_formula.gp} / \texttt{.out} --- reproduction of the heuristic
threshold \eqref{eq:N0} on all fibers (used for context, not for closure).
\end{itemize}

\section*{Acknowledgements}
The author thanks the maintainers of PARI/GP.

\begin{thebibliography}{9}

\bibitem{ingram07}
P.~Ingram,
\emph{Elliptic divisibility sequences over certain curves},
J. Number Theory \textbf{123} (2007), no.~2, 473--486.

\bibitem{ingram-silverman}
P.~Ingram and J.~H.~Silverman,
\emph{Uniform estimates for primitive divisors in elliptic divisibility
sequences},
in: Number Theory, Analysis and Geometry (in memory of Serge Lang),
Springer, New York, 2012, pp.~243--271.

\bibitem{matson}
R.~D.~Matson,
\emph{Results of a computer search for a perfect cuboid},
unpublished manuscript, 2009;
\url{https://www.unsolvedproblems.org/}.

\bibitem{mazur}
B.~Mazur,
\emph{Rational isogenies of prime degree},
Invent. Math. \textbf{44} (1978), 129--162.

\bibitem{pari}
The PARI~Group,
\emph{PARI/GP version 2.15.4},
Univ. Bordeaux, 2023,
\url{https://pari.math.u-bordeaux.fr/}.

\bibitem{peschmann1}
R.~Peschmann,
\emph{A torsion-intersection proof of perfect-cuboid nonexistence on $1{,}072$
explicit master-tuple fibers},
preprint, arXiv:2604.28072 (2026).

\bibitem{peschmann2}
R.~Peschmann,
\emph{Quartic reductions and elliptic obstructions for perfect Euler bricks},
preprint, arXiv:2604.09328 (2026).

\bibitem{silverman88}
J.~H.~Silverman,
\emph{Wieferich's criterion and the $abc$-conjecture},
J. Number Theory \textbf{30} (1988), no.~2, 226--237.

\bibitem{silverman-eds}
J.~H.~Silverman,
\emph{$p$-adic properties of division polynomials and elliptic divisibility
sequences},
Math. Ann. \textbf{332} (2005), no.~2, 443--471.

\bibitem{verzobio}
M.~Verzobio,
\emph{Some effectivity results for primitive divisors of elliptic divisibility
sequences},
Pacific J. Math. \textbf{325} (2023), no.~2, 331--351; arXiv:2001.02987.

\bibitem{yoshida}
T.~Yoshida,
\emph{The relationship between face cuboids and elliptic curves},
preprint, arXiv:2407.09825 (2024).

\end{thebibliography}

\end{document}
